\section{角度の倍増}\label{節：角度の倍増}
\subsection{角度を倍にする場合}

次は
\[
 a_1=\frac{\sqrt3}{2},\qquad a_{n+1}=2a_n^2-1
\]
を考える.
倍角公式$\cos2\theta=2\cos^2\theta-1$と,
$a_1=\cos\frac{\pi}{6}$を使えば,
\[
 a_n=\cos\left(2^{n-1}\frac{\pi}{6}\right)
\]
と予想できる.
初項に合い,\,この式を更新の右辺に代入すると次の項の式になるので,\,帰納法で正しい.

角度は$\theta_{n+1}=2\theta_n$と選べる.
複素数では
\[
 z_{n+1}=z_n^2,\qquad z_n=z_1^{\,2^{n-1}}
\]
であり,$a_n=\operatorname{Re}z_n$である.
\NewFigDoubling

これは一定角の回転ではない.
現在の角度$\theta_n$をもう一度足すため,\,毎回の回転角が変わる.
さらに$z$と$-z$は同じ$z^2$へ移るので,\,円周全体の写像として一対一ではない.

余弦の値から角度は一意に決まらない.
そのため$a_{n+1}=2a_n^2-1$から,\,選び方によらず$\theta_{n+1}=2\theta_n$が強制されると読むのは正確でない.
整合する角度の列を一つ選び,\,そこから作った候補が元の式を満たすことを示している.

\subsection{正接の倍角と定義域}

傾きを記録する倍角型は
\[
 a_{n+1}=\frac{2a_n}{1-a_n^2}
\]
である.
$a_n=\tan\theta_n$と表せれば,$\theta_{n+1}=2\theta_n$と選べる.
ただし$a_n=\pm1$では次の傾きが縦向きになり,\,実数の値としては定義できない.
円周上の動きが続いても,\,傾きの列が続くとは限らない.

\noindent\textbf{【例題】}\par\nobreak
\[
 a_1=2-\sqrt3,\qquad
 a_{n+1}=\frac{2a_n}{1-a_n^2}
\]
の一般項を求め,\,全項が定義されることを示せ.

\noindent\textbf{【方針】}\par\nobreak
$2-\sqrt3=\tan\frac{\pi}{12}$です. 角度を倍にし,\,途中で縦向きにならないか確認します.

\noindent\textbf{【解答】}\par\nobreak
\[
 a_n=\tan\left(2^{n-1}\frac{\pi}{12}\right).
\]
\noindent\textbf{【解法】}\par\nobreak
初項と倍角公式から候補を得る.
$n=1$では正接は定義できる.
$n\ge2$でその角度が$\frac{\pi}{2}+\ell\pi$に一致するとすれば
\[
 2^{n-2}=3+6\ell
\]
となるが,2の累乗は3の倍数でないので不可能である.
従って全項の正接が定義される.
また次の角度も縦向きにならないため,\,更新の分母$1-a_n^2$も非零である.
帰納法により一般項が成立する.

なお,$\Phi$は\sectionref{節：実軸と円を結ぶ座標}の実軸と円の対応であり,
\[
 \Phi\left(\frac{2a}{1-a^2}\right)=\Phi(a)^2
\]
となる.
正接の倍角型も,\,円周上の二乗を別の座標で記録していた.
一次分数型の一定角回転とは,\,同じ単位円を使っても更新の規則が違う.


\begin{bookpoint}
一定角を足す動きでは,複素数は一定数を掛けられる.
角度を倍にする動きでは,複素数は累乗される.
同じ三角関数が使えても,二つの更新を区別して考える.
\end{bookpoint}
